%% kapa: exam=1 round=25 year=2014 period=1 form=A subject=civil_law subject_ko=민법 seq=34 number=34
%% source: Q-Net 감정평가사 기출문제 (cst00309) artlSeq=5101207 file=extracted/1008/5101207/2014년 제25회 감정평가사 1차시험 기출문제/1교시(A형).hwp
%% conversion: kordoc 4.24.51 -> tools/gen1.py (auto); review=proofread-sonnet; flags=none
\begin{question}{34}
지역권에 관한 설명으로 옳지 않은 것은?
\begin{choices}
  \item 지역권은 소멸시효의 대상이 될 수 없다.
  \item 공유자의 1인이 지역권을 취득한 때에는 다른 공유자도 이를 취득한다.
  \item 승역지 소유자는 지역권의 행사를 방해하지 않는 범위 내에서 지역권자가 지역권의 행사를 위하여 승역지에 설치한 공작물을 사용할 수 있다.
  \item 지역권은 요역지 소유권의 상속에 의하여 취득될 수 있다.
  \item 지역권은 다른 약정이 없는 한, 요역지 소유권에 부종하여 이전한다.
\end{choices}
\end{question}
\answer{1}
